% conclusions.tex — §8 (merged discussion + conclusions)

\section{Discussion and Conclusions}
\label{sec:conclusions}

We have applied Singular Value Decomposition to DESI DR2 BAO measurements and three Type~Ia supernova compilations (Union3, Pantheon+, DES-Dovekie), all analyzed with a common CMB baseline from ACT DR6. Our main findings are:

\begin{enumerate}

\item \textbf{Universal \czero.} All low-$z$ distance probes measure the same leading SVD mode \czero which we loosely call \omh (inner products $> 0.996$; $R^2 > 0.99$). A single number captures what these four experiments collectively say about cosmology.

\item \textbf{The \wowa preference is primarily \czero.} The DESI \wowa preference is dominated by \czero; the BAO dark energy direction (\cone) shows only $+1.2\sigma$ tension with \LCDM. At the BAO \cone pivot ($z = 0.46$), $w = -0.94 \pm 0.052$ ($+1.2\sigma$)---consistent with a cosmological constant. SN \cone tensions are also mild ($-1.5\sigma$ to $+0.2\sigma$). BAO alone prefers lower \omh ($+2.2\sigma$), while all three SN datasets are mild or null ($-0.8$ to $-1.7\sigma$), none confirming the BAO signal.

\item \textbf{Curvature coherence.} Spatial curvature is the only other extension that opens a new BAO direction although barely; \czero and \cone independently point to positive spatial curvature an interesting consistency check, altough the errorbar associated with the \cone direction is very large. However, SN are blind to curvature ($\sigres < 0.1$), so cross-probe confirmation is impossible.

\item \textbf{No other extension is distinguishable.} No other tested extension (EDE, $A_\mathrm{lens}$, $B_\mathrm{PMF}$) opens a measurable new direction ($\sigres < 0.3$). Low-$z$ distance data cannot distinguish among them.
 
\item \textbf{CMB lensing contribution.} The CMB \omh constraint has a lensing component: $A_\mathrm{lens}$ reduces BAO tension from $2.2\sigma$ to $1.1\sigma$ by both broadening and shifting the \omh posterior. This is the same physics as shifting the optical depth $\tau$ in the CMB. 

\end{enumerate}

The SVD framework tells us \emph{where} the tension is (\czero, i.e., \omh) but not \emph{what} causes it. This is a fundamental limitation of distance measurements, not a shortcoming of the method: the SVD transparently identifies the boundary between what low-$z$ probes can and cannot constrain.

%Our \czero corresponds to the Model-Embedding Distance Indicator (MEDI) of \citet{Weiner2026_MEDI} (inner product 0.97). The MEDI framework provides the analytic explanation for \emph{why} \czero has suppressed dark energy sensitivity---dark energy corrections enter at $\sim 1/14$ of the \omh effect---and identifies specific high-$z$ resolution mechanisms. Our SVD complements this by adding empirical cross-probe universality, the BAO--SN sign pattern pointing to systematics, and the curvature coherence test from two independent BAO directions. The two approaches reach the same conclusion: the DESI dark energy signal is an \omh calibration discrepancy.

All current low-redshift distance data, whether from BAO or supernovae, are effectively one-dimensional in their cosmological information content. That single dimension is well described by \omh.
